\section{Choose the accepted image subset}
The exercise accepts static, noninterlaced RGB images with eight bits per channel
(PNG color type 2, bit depth 8), dimensions from 1 through 512, and at most 1 MiB of
file input. Compression and filter methods are both zero. Other valid PNG formats,
including indexed color, grayscale, alpha channels and interlacing, are outside this
exercise. Reject them before reserving row storage or decoding compressed samples.
The bitmap output preserves sample numbers; it does not perform color management.

Require IHDR first, one uninterrupted run of IDAT chunks, and an empty IEND at the
exact end of the supplied file. An optional PLTE may occur once before IDAT; its
payload must contain 1 through 256 RGB triples. Truecolor pixels do not index that
palette. Check every chunk CRC, including chunks whose payload you skip. Validate
that all four type bytes are ASCII letters and that the reserved third letter is
uppercase. Reject unknown critical chunks. Skip otherwise permitted ancillary chunks
without interpreting their data. An ancillary chunk after IDAT closes the IDAT run.

This teaching subset also rejects transparency and animation control/data chunks
\texttt{tRNS}, \texttt{acTL}, \texttt{fcTL} and \texttt{fdAT}. Silently ignoring them
would make the displayed result misleading. These are local restrictions, not claims
that such PNG files are malformed. Likewise, skipping color metadata is not equivalent
to implementing its color interpretation. The later APNG chapter will replace the
static-image restriction with explicit frame semantics.
